% Q04_solutions.tex — REFERENCE SOLUTIONS to paper/Q04_problems.tex.
% Author-only and read-denied in study mode (Read(./paper/solutions/**)).
% Built only by `make -C paper paper-solutions` (scripts/ci.sh, or the
% author's linked worktree); `make paper` never touches this directory.
%
% Every verifier result quoted below was captured on the course pins (Dafny
% 4.11.0, Verus 0.2026.09.27.3cf1832) through scripts/vcheck and scripts/vaudit,
% run on instructor/scratch/probe.<ext> in a scratch course root, one solver job
% at a time under a 16G cap, by the DV3a kits phase's kit lane for Q04 on
% 2026-10-06. The probe files and their outputs are archived under
% ~/.cache/dafny-verus-learn-0/author-probes/dv3k/q04-kit/ (README there).
% B3's fields are graded against the learner's own cheatsheets.md; this file
% records which lines of the pinned range each field draws on and never copies
% the card.
\documentclass[11pt]{article}
\input{preamble}

\title{Q04 --- Proof craft, loops and libraries \\ \large Reference solutions (author-only)}
\author{dafny-verus-learn-0}
\date{}

\begin{document}
\maketitle

\section*{How these are graded}
B1 is informal-proof-first (the five proof dimensions; its blueprint is
re-scored by \verb+--post+ against the learner's verified H4d and H4v). B2 is a
concept answer (\texttt{correctness}: the grader tries to counterexample it,
with the rubric's Probe B). B3 is a paper card (\texttt{card-fidelity}, per
field). The solutions below give the answers; the rubric never states them.
The structure of B1 matches the formal references of H4d and H4v (spec
section~4.4's author rule).

% ---------------------------------------------------------------------------
\section*{B1 (informal-proof-first)}

\textbf{(a) Statement.} For all limb sequences $a$, $b$ with
$\mathit{Valid}(a)$ and $\mathit{Valid}(b)$: $-1 \leq \mathit{Cmp}(a, b) \leq 1$,
$\mathit{Cmp}(a, b) < 0 \iff \mathit{ToNat}(a) < \mathit{ToNat}(b)$, and
$\mathit{Cmp}(a, b) = 0 \iff \mathit{ToNat}(a) = \mathit{ToNat}(b)$.

\textbf{Induction on} $|a| + |b|$, the measure of \dfy{Cmp} itself and the
lemma's shipped \dfy{decreases}.

\textbf{IH (as a formula).} For all valid $a'$, $b'$ with
$|a'| + |b'| < |a| + |b|$: the three claims hold for $(a', b')$. It is used
once, at $(\mathit{Tail}(a), \mathit{Tail}(b))$, and only when
$|a| + |b| > 0$: then at least one sequence is non-empty, \dfy{Tail} removes
one limb from it and leaves an empty sequence empty, so
$|\mathit{Tail}(a)| + |\mathit{Tail}(b)| < |a| + |b|$.

\textbf{Base} ($|a| = |b| = 0$). $\mathit{Cmp}(a, b) = 0$ and both values are
$0$: the range holds, $0 < 0$ is false on both sides, and $0 = 0$ on both.

\textbf{Step} ($|a| + |b| > 0$). Write $h_a = \mathit{Head}(a)$,
$h_b = \mathit{Head}(b)$, $t_a = \mathit{ToNat}(\mathit{Tail}(a))$,
$t_b = \mathit{ToNat}(\mathit{Tail}(b))$ and
$c = \mathit{Cmp}(\mathit{Tail}(a), \mathit{Tail}(b))$. First the facts:
\begin{enumerate}
\item L2's \dfy{ValidTail(a)} and \dfy{ValidTail(b)}: the tails are valid and
  $h_a, h_b < \mathit{BASE}$ (for an empty sequence too: its head reads as 0);
\item the IH at the tails, whose \dfy{requires} item~1 supplies:
  $c \in \{-1, 0, 1\}$, $c < 0 \iff t_a < t_b$, $c = 0 \iff t_a = t_b$;
\item L2's \dfy{ToNatHeadTail(a)} and \dfy{ToNatHeadTail(b)}:
  $\mathit{ToNat}(a) = h_a + \mathit{BASE} \cdot t_a$ and
  $\mathit{ToNat}(b) = h_b + \mathit{BASE} \cdot t_b$.
\end{enumerate}
\emph{The tails decide} ($c \neq 0$, so $\mathit{Cmp}(a, b) = c$). If $c < 0$
then $t_a < t_b$, so $t_a + 1 \leq t_b$ over the naturals, and
\[
  \mathit{ToNat}(a) = h_a + \mathit{BASE} \cdot t_a
  < \mathit{BASE} \cdot (t_a + 1) \leq \mathit{BASE} \cdot t_b
  \leq h_b + \mathit{BASE} \cdot t_b = \mathit{ToNat}(b),
\]
the strict step by $h_a < \mathit{BASE}$: one unit of the tails is worth
$\mathit{BASE}$, more than any head. So $\mathit{Cmp}(a, b) = -1 < 0$ and
$\mathit{ToNat}(a) < \mathit{ToNat}(b)$; neither side of the third claim
holds. If $c > 0$ then $c = 1$ and $t_b < t_a$; the same chain with $a$ and
$b$ exchanged gives $\mathit{ToNat}(b) < \mathit{ToNat}(a)$, and all three
claims hold with both sides of the second and third false.

\emph{The heads decide} ($c = 0$, so $t_a = t_b$ and
$\mathit{ToNat}(a) - \mathit{ToNat}(b) = h_a - h_b$). \dfy{Cmp} returns $-1$,
$1$ or $0$ as $h_a < h_b$, $h_a > h_b$ or $h_a = h_b$, and the values compare
in the same way.

\textbf{Helper lemmas (the blueprint).} L2's \dfy{ValidTail} (on $a$ and on
$b$) and \dfy{ToNatHeadTail} (on $a$ and on $b$), and the recursive call at the
tails. The dominance step ($h_a < \mathit{BASE}$ and $t_a < t_b$ give
$h_a + \mathit{BASE} \cdot t_a < h_b + \mathit{BASE} \cdot t_b$) is linear
arithmetic over the constant \dfy{BASE}, which the verifier does unaided;
naming it as a new lemma with exactly those hypotheses is equally correct
(captured: the gold fixture's \dfy{Dominance}, empty body,
\texttt{verified --- 23 verified, 0 holes, 0 errors}).

\textbf{Why the validity hypothesis is needed.} Without $h_a < \mathit{BASE}$
the tails need not decide: $a = [\mathit{BASE}]$, $b = [0, 1]$ have equal
values ($\mathit{BASE}$) and $\mathit{Cmp}(a, b) = -1$ (captured: the rubric's
Probe A$''$ with \dfy{expect (Cmp(a, b) == 0) == (ToNat(a) == ToNat(b))} at
these sequences, \texttt{FAILED}).

\textbf{Gaps.} None.

\textbf{The formal reference it matches} (H4d's twin):
\begin{dafnycode}
lemma CmpCorrect(a: seq<nat>, b: seq<nat>)
  requires Valid(a) && Valid(b)
  ensures -1 <= Cmp(a, b) <= 1
          && (Cmp(a, b) < 0 <==> ToNat(a) < ToNat(b))
          && (Cmp(a, b) == 0 <==> ToNat(a) == ToNat(b))
  decreases |a| + |b|
{
  if |a| != 0 || |b| != 0 {
    ValidTail(a);
    ValidTail(b);
    CmpCorrect(Tail(a), Tail(b));
    ToNatHeadTail(a);
    ToNatHeadTail(b);
  }
}
\end{dafnycode}
The base case needs no line, and the case split on $c$ is \dfy{Cmp}'s own
unfolding. This blueprint spliced into the rubric's Probe A (the target under
the name \dfy{CmpAgrees}) gave \texttt{verified --- 31 verified, 0 holes,
0 errors} under \texttt{scripts/vcheck --module Q04} (the count includes the
given layers' tasks).

\textbf{(b) Verus.}
\begin{itemize}
\item \textbf{The missing clause:} \vrs{forall|j: int| i <= j < n ==> a@[j] == b@[j]},
  the limbs from index $i$ to the top agree.
\item \textbf{On entry} $i = n$ and the range is empty. \textbf{One pass}
  starts where the guard holds, so $i > 0$ and $a[i - 1] = b[i - 1]$, and
  ends with $i$ one smaller: the range gains exactly index $i - 1$ of the start
  of the pass, where the guard gave equality.
\item \textbf{After the loop}, with equal lengths ($n$ = both lengths):
  if $i = 0$, the clause covers every index, so the views are equal
  (\vrs{a@ =~= b@}) and so are the values: \vrs{Equal}. Otherwise the guard
  failed with $i > 0$, so $a[i - 1] \neq b[i - 1]$. When $a[i - 1] < b[i - 1]$,
  \vrs{lemma_lt_at(a@, b@, i - 1)}'s four \vrs{requires} hold: equal lengths
  (the second and third invariant clauses), $0 \leq i - 1 < |a|$ ($i \leq n$),
  $a[i - 1] < b[i - 1]$ (the branch), and agreement at every index above
  $i - 1$, which is the missing clause. So $a$'s value is the smaller:
  \vrs{Less}. The other branch has $b[i - 1] < a[i - 1]$ and calls the lemma
  with the views exchanged: \vrs{Greater}. Each outcome meets all three
  \vrs{ensures}, since the value order is total.
\item \textbf{The length cases.} For $n < m$: $b$ is canonical and non-empty
  ($m > n \geq 0$), so L3's \vrs{lemma_seq_to_nat_lower} gives
  $\mathit{weight}(m - 1) \leq \mathit{value}(b)$; L3's
  \vrs{lemma_seq_to_nat_upper} gives $\mathit{value}(a) <
  \mathit{weight}(n)$; and L3's \vrs{lemma_weight_mono} at $n \leq m - 1$ gives
  $\mathit{weight}(n) \leq \mathit{weight}(m - 1)$. Chained:
  $\mathit{value}(a) < \mathit{weight}(n) \leq \mathit{weight}(m - 1) \leq
  \mathit{value}(b)$: \vrs{Less}. For $n > m$ the same with $a$ and $b$
  exchanged: \vrs{Greater}. Canonicity is what makes length decide: $[5, 0]$
  is longer than $[7]$ and smaller.
\item \textbf{How the formal proof gets the three facts.} They are the members
  of L3's group \vrs{group_bignat_bounds}; H4v's twin brings them into the
  whole body with one \vrs{broadcast use group_bignat_bounds;} line at its top,
  and their triggers fire on the values and weights the two length cases
  mention. A blueprint that names the three lemma calls instead is equally
  correct: the rubric's Probe A$'$ with this clause and the named calls in the
  two length branches gave \texttt{verified --- 2 verified, 0 holes, 0 errors}.
\end{itemize}
The formal reference's two learner regions (H4v's twin): the first line of the
body, and the invariant's last clause.
\begin{veruscode}
    broadcast use group_bignat_bounds;
    // ...
        invariant
            i <= n,
            n == a@.len(),
            n == b@.len(),
            forall|j: int| i <= j < n ==> a@[j] == b@[j],
\end{veruscode}
An equivalent formula is not always formalizable as written: the clause
written as \vrs{a@.subrange(i as int, n as int) =~= b@.subrange(i as int, n as int)}
is true in every state the loop reaches, yet Probe A$'$ with it fails four
ways (the invariant at the end of the body, both \vrs{lemma_lt_at}
preconditions and the equality assertion), because nothing then states the
per-index agreement \vrs{lemma_lt_at} quantifies over. That is a gap (3) for
\texttt{blueprint-quality}, never a flaw.

% ---------------------------------------------------------------------------
\section*{B2 (concept)}

\textbf{(a)} The outcomes are the same in both languages (each captured
with \texttt{scripts/vcheck --module Q04 --gate} and \texttt{scripts/vaudit}
on the probe):
\begin{center}
\begin{tabular}{llll}
body & \texttt{vcheck --gate} & \texttt{vaudit} & ``only library calls'' \\
\hline
{[1]} call alone & passes & 0 findings & kept \\
{[2]} call inside an assertion & passes & 0 findings & broken \\
{[3]} an assumption & fails: the assumption is a hole & 1 finding & broken \\
\end{tabular}
\end{center}
For [3] the gate prints a hole line for the assumption in each language
(Dafny: \texttt{hole assume ProbeB2.ScaleLe}, with \texttt{warnings-fail-gate};
Verus: \texttt{hole assume lemma\_scale\_le}), and the audit's finding is
\texttt{audit-assume} in Dafny and \texttt{hole-assume} in Verus. [2] is the case no instrument sees: the
assertion states the postcondition, so the proof is correct, green and
audit-clean, and the contract's acceptance line forbids exactly this (``no
\dfy{assert}''). With nonlinear arithmetic left on, the Dafny empty body also
verifies (captured), which is why the race's Dafny file turns it off.

\textbf{(b)} Any body that passes both and is not a bare sequence of library
calls. The reference: a helper lemma of one's own with the goal's statement,
whose body is the library call, and the goal's body calling the helper
(captured: Dafny \texttt{verified --- 32 verified, 0 holes, 0 errors} and
\texttt{0 finding(s)}; Verus \texttt{verified --- 2 verified, 0 holes, 0
errors} and \texttt{0 finding(s)}). A \dfy{calc} or \vrs{calc!} chain with the
call as a hint is another. The breach shows in \texttt{git diff} (an added
declaration outside the bodies, or a statement inside one that is not a
call), which the contract's diff and body lines ask the learner to read, and
in the transcript; the gate, the audit and the hole count all stay clean.
A claim that the audit or the gate reports [2] or the fourth body is refuted
by running them (the rubric's Probe B).

% ---------------------------------------------------------------------------
\section*{B3 (paper card 02.2): reference fields and their lines}

Pinned range: \texttt{cheatsheets.md} lines 614--638 (computed 2026-10-06),
\[ \texttt{sha256 5a0ff58db586d31eb8bde23fb09f5e1c8f4d2a558b339bf970050c8aba407e6a}. \]
A field is correct when it agrees with the lines named; wording is free. Each
entry: the reference content (paraphrased), then the lines it draws on.

\begin{description}
\item[Task shape.] Input: a Rust function with its Verus specification
  (\vrs{requires}/\vrs{ensures}). Output: the same function with proof
  annotations added (loop invariants, assertions, \vrs{proof fn} lemmas) until
  Verus verifies it; GPT-4o agents in three phases (generate and merge
  invariants, refine, repair errors by type). \emph{Lines 616, 618, 622--624.}
\item[Oracle.] Verus's verdict, read as a score: $V$, the number of verified
  functions (each loop counting as one), and $E$, the number of errors, rank
  and accept candidates; Lynette's check that a candidate changed neither code
  nor specification and added no \vrs{assume} gates them.
  \emph{Lines 618, 621, 623--624.}
\item[What is trusted.] Verus and its solver; Lynette, the authors' AST tool;
  the specifications themselves (the limitations say so); the hand-crafted
  prompts shape what is tried. Whether other escape hatches (an external body,
  an \vrs{admit}) are rejected is not stated. \emph{Lines 618, 636.}
\item[Cheat surface.] Closed by Lynette: editing the code, editing the
  specification, adding \vrs{assume}. Open or unstated: a weak or wrong
  specification is trusted, so nothing catches it; escape hatches other than
  \vrs{assume}; redundant invariants are harmless noise, not a cheat.
  \emph{Lines 618, 636.}
\item[Metric.] Tasks proved out of Verus-Bench's 150: 137 for AutoVerus (best
  of 3 runs; 122 in the first run), 67 for a GPT-4o prompting baseline;
  ablations in task counts (a generic debugger 17/37 against error-specific
  agents 37/37; 26, 26, 23 and 18 of a 30-task sample by model).
  \emph{Lines 616, 618, 628--634.}
\item[Budget.] The baseline had 10 minutes per task; to reach the baseline's
  67 tasks AutoVerus needed 22 s and 2 calls, the baseline 458 s and 59 calls;
  inside the pipeline, 5 samples in phase 1 and up to 10 repair iterations in
  phase 3. No token, money or per-task time figure for the 137.
  \emph{Lines 622, 624, 631--632.}
\item[Variance.] Three runs; the headline is the best of them, and the first
  run alone proved 122, so the spread (122 to 137) is the only variance
  signal; no mean, interval or seed count; sampling at temperature 1.0.
  \emph{Lines 622, 630, 636.}
\item[Artifact.] Verus-Bench, 150 tasks from four sources (MBPP 78, Diffy 38,
  Misc 23, CloverBench 11), introduced here and reused by later work (a
  139-task variant in SAFE; AlphaVerus compares against AutoVerus). Whether
  the code or prompts are released is not stated. \emph{Lines 616, 618, 638.}
\end{description}
A field that says the range is silent where it is silent (a released
artifact; a cost for the headline; an interval) is correct. A field that
contradicts a line (for example 137 as a single run, or the 22 s and 2 calls
given to the baseline) scores 2, quoting that line.

\end{document}
